\section{Turing Machine --- Build --- Window Run}
%%% generated by the blueprint pipeline from TCSlib.Complexity.TuringMachine.Build.WindowRun
\subsection*{Turing Machine --- Build --- Window Run: Introduction}

Declarations of the Lean module \texttt{TCSlib.Complexity.TuringMachine.Build.WindowRun}.

\subsection*{Turing Machine --- Build --- Window Run: Results}

\begin{definition}[Window exit outcomes]
\lean{Turing.WindowExit}
\label{Turing.WindowExit}
\leanok
A window exit has one of three possible outcomes: the simulated machine has halted, its
next simulated step leaves the window, or the clock is exhausted.
\end{definition}

\begin{definition}[Window-limited machine run]
\lean{Turing.windowRun}
\label{Turing.windowRun}
\leanok
\uses{Turing.Cfg, Turing.Cfg.InWindow, Turing.MultiTapeTM, Turing.MultiTapeTM.outputSymbol, Turing.MultiTapeTM.step, Turing.WindowExit}
Let $M$ be a deterministic Turing machine with $k$ work tapes and Boolean symbols, let
$s$ be a window radius, and let $\mathrm{upd}:T\times\mathrm{Option}(\mathrm{Bool})\to
T$ update a tag. Define the run for a clock value $N$ and configuration $c$ with tag
$\tau$ as follows: at $N=0$, return the exhausted exit together with $c$ and $\tau$. At
positive clock, if $c$ is halted, return the halted exit with $c$ and $\tau$; otherwise,
if one step from $c$ lies within the window of radius $s$, continue for one fewer clock
step from that successor configuration, updating the tag using the symbol emitted by the
step from $c$. If the successor lies outside the window, return the left exit with the
current configuration $c$ and tag $\tau$.
\end{definition}

\begin{theorem}[Window run through a halted configuration]
\lean{Turing.windowRun_halted_of_run}
\label{Turing.windowRun_halted_of_run}
\leanok
\uses{Turing.Cfg, Turing.Cfg.InWindow, Turing.MultiTapeTM, Turing.MultiTapeTM.runFrom, Turing.MultiTapeTM.runFrom_of_halt, Turing.MultiTapeTM.runFrom_succ_eq_step, Turing.MultiTapeTM.step, Turing.MultiTapeTM.step_output, Turing.windowRun}
\difficulty{5}
Let $M$ be a Boolean-input Turing machine with $k$ work tapes, let $s$ be a
natural-number window radius, and let $\mathrm{upd}$ be an update rule on an accumulator
type $T$ and an optional Boolean. Suppose $\sigma$ assigns an accumulator to every
Boolean word and satisfies $\sigma(u\mathbin{++}b)=\mathrm{upd}(\sigma(u),b)$ for every
word $u$ and optional Boolean $b$, where appending an absent Boolean leaves $u$
unchanged. For every clock $N$, time $t$, and initial configuration $c$, if $t<N$, the
run is halted at time $t$, and all configurations through time $t$ lie in the radius-$s$
window, then the window run with initial accumulator $\sigma$ of the initial output
halts and returns the configuration at time $t$ together with accumulator $\sigma$ of
its output.
\end{theorem}

\begin{theorem}[Windowed run characterization]
\lean{Turing.windowRun_spec}
\label{Turing.windowRun_spec}
\leanok
\uses{Turing.Cfg, Turing.Cfg.InWindow, Turing.MultiTapeTM, Turing.MultiTapeTM.outputSymbol, Turing.MultiTapeTM.runFrom, Turing.MultiTapeTM.runFrom_succ_eq_step, Turing.MultiTapeTM.runFrom_zero, Turing.MultiTapeTM.step, Turing.windowRun}
\difficulty{6}
Let $M$ be a multi-tape Turing machine with Boolean symbols, let $s,N\in\mathbb{N}$, let
$\mathrm{upd}$ update an auxiliary value from an optional Boolean symbol, and let $c$ be
a configuration and $\tau$ an auxiliary value, with $c$ in the window of radius $s$.
Then there is an $r\le N$ such that the configuration returned by the windowed run is
exactly the configuration reached by $M$ after $r$ steps; every configuration up to and
including step $r$ is in the window, and every configuration before step $r$ is
nonhalting. If the run reports halting, then $r<N$ and the configuration at step $r$ is
halted; if it reports leaving the window, then $r<N$, the configuration at step $r$ is
nonhalting, and its next configuration is outside the window; if it reports exhaustion,
then $r=N$.
\end{theorem}

\begin{theorem}[Accumulator tracks the output word]
\lean{Turing.windowRun_tag}
\label{Turing.windowRun_tag}
\leanok
\uses{Turing.Cfg, Turing.Cfg.InWindow, Turing.MultiTapeTM, Turing.MultiTapeTM.step, Turing.MultiTapeTM.step_output, Turing.windowRun}
\difficulty{4}
Let $M$ be a multi-tape Turing machine with Boolean symbols, let $s\in\mathbb{N}$, and
let $\mathrm{upd}$ update an auxiliary value from an optional Boolean symbol. Suppose
$\sigma$ assigns an auxiliary value to each Boolean word and satisfies
$\sigma(u\mathbin{+\!\!+}\mathrm{toList}(b))=\mathrm{upd}(\sigma(u),b)$ for every Boolean word $u$ and
optional Boolean symbol $b$, where $\mathrm{toList}(b)$ is the empty word if $b$ is absent and the one-symbol word containing $b$ if it is present.
For every clock bound $N$ and configuration $c$, the auxiliary value returned by the
windowed run started with $\sigma$ of the initial output word equals $\sigma$ of the
output word in the returned configuration.
\end{theorem}

\begin{theorem}[Uniqueness of non-exhausted window runs]
\lean{Turing.windowRun_stop_unique}
\label{Turing.windowRun_stop_unique}
\leanok
\uses{Turing.Cfg, Turing.Cfg.InWindow, Turing.MultiTapeTM, Turing.MultiTapeTM.step, Turing.windowRun}
\difficulty{4}
Let $M$ be a multi-tape Turing machine with Boolean symbols, let $s\in\mathbb{N}$, let
$\mathrm{upd}$ update an auxiliary value from an optional Boolean symbol, and let
$N,N'\in\mathbb{N}$. For every configuration $c$ and auxiliary value $\tau$, if the
windowed runs with clock bounds $N$ and $N'$ both report an outcome other than
exhaustion, their returned outcome, configuration, and auxiliary value are equal.
\end{theorem}

\begin{theorem}[Exhaustion of a live window run]
\lean{Turing.windowRun_exhausted_of_live}
\label{Turing.windowRun_exhausted_of_live}
\leanok
\uses{Turing.Cfg, Turing.Cfg.InWindow, Turing.MultiTapeTM, Turing.MultiTapeTM.outputSymbol, Turing.MultiTapeTM.runFrom, Turing.MultiTapeTM.runFrom_succ_eq_step, Turing.MultiTapeTM.step, Turing.windowRun}
\difficulty{4}
Let $M$ be a Boolean-input Turing machine with $k$ work tapes, let $s$ be a
natural-number window radius, and let $\mathrm{upd}$ be an update rule on an accumulator
type $T$ and an optional Boolean. For every natural number $N$, initial configuration
$c$, and accumulator $\tau$, suppose the machine is not halted at any time $u<N$ and its
configuration is in the window of radius $s$ at every time $u\le N$. Then the window run
with clock $N$ returns the exhausted outcome and its configuration component is exactly
the configuration reached by $M$ after $N$ steps from $c$.
\end{theorem}

\begin{theorem}[Halting on the final transition]
\lean{Turing.windowRun_halts_on_last_transition}
\label{Turing.windowRun_halts_on_last_transition}
\leanok
\uses{Turing.Cfg, Turing.Cfg.InWindow, Turing.MultiTapeTM, Turing.MultiTapeTM.outputSymbol, Turing.MultiTapeTM.runFrom, Turing.MultiTapeTM.runFrom_succ_eq_step', Turing.MultiTapeTM.step_output, Turing.windowRun, Turing.windowRun_halted_of_run}
\difficulty{3}
Let $M$ be a Boolean-input Turing machine with $k$ work tapes, let $s$ be a
natural-number window radius, and let $\mathrm{upd}$ be an update rule on an accumulator
type $T$ and an optional Boolean. Suppose $\sigma$ assigns an accumulator to every
Boolean word and satisfies $\sigma(u\mathbin{++}b)=\mathrm{upd}(\sigma(u),b)$ for every
word $u$ and optional Boolean $b$, where appending an absent Boolean leaves $u$
unchanged. If, for a time $t$, initial configuration $c$, and Boolean $b$, the run is
halted at time $t+1$, the transition from time $t$ emits $b$, and all configurations
through time $t+1$ lie in the radius-$s$ window, then the window run with clock $t+2$
halts and returns the configuration at time $t+1$ together with the accumulator obtained
by applying $\sigma$ to the time-$t$ output followed by $b$.
\end{theorem}

\begin{definition}[Window-loop phases]
\lean{Turing.wl_Phase}
\label{Turing.wl_Phase}
\leanok
Define a four-element phase type with constructors for the round anchor, the test phase,
the step phase, and the stop phase.
\end{definition}

\begin{definition}[Window-loop control states]
\lean{Turing.wl_Ctrl}
\label{Turing.wl_Ctrl}
\leanok
\uses{Turing.WindowExit, Turing.wl_Phase}
For types $P_t$ and $P_s$, define the control-state type indexed by the four window-loop
phases: it is the unit type at the anchor phase, $P_t$ at the test phase, $P_s$ at the
step phase, and the three-valued window-exit type at the stop phase.
\end{definition}

\begin{definition}[Window-loop record tape]
\lean{Turing.wl_recordTape}
\label{Turing.wl_recordTape}
\leanok
For each natural number $K$, define an embedding of the one-element index type into the
$(K+1)$-element index type by sending its unique element to index zero.
\end{definition}

\begin{definition}[Window-loop transition machines]
\lean{Turing.wl_rows}
\label{Turing.wl_rows}
\leanok
\uses{Turing.MultiTapeTM, Turing.RecSym, Turing.embedEmitTM, Turing.wl_Ctrl, Turing.wl_Phase, Turing.wl_recordTape}
Given a one-tape test machine $R_t$ and a $(K+1)$-tape step machine $R_s$, define a
machine for each window-loop phase, all using $K+1$ work tapes and the record-symbol
alphabet. At the anchor phase it remains in its unique state without moving or writing;
at the test phase it runs $R_t$ embedded on tape zero; at the step phase it is $R_s$;
and at the stop phase each window-exit state remains fixed without moving or writing.
\end{definition}

\begin{definition}[Window-loop phase transitions]
\lean{Turing.wl_red}
\label{Turing.wl_red}
\leanok
\uses{Turing.WindowExit, Turing.wl_Ctrl, Turing.wl_Phase}
Let $P_t$ and $P_s$ be types with decidable equality, and choose test states
$t_e,t_y,t_n\in P_t$ and step states $s_e,s_y,s_n\in P_s$. The phase-dependent
transition sends every anchor control to test state $t_e$; in the test phase it sends
$t_y$ to the halted stop, $t_n$ to step state $s_e$, and every other control to no
successor. In the step phase it sends $s_y$ back to the anchor, $s_n$ to the left stop,
and every other control to no successor; all stop controls have no successor.
\end{definition}

\begin{definition}[Window-loop machine construction]
\lean{Turing.wl_loopTM}
\label{Turing.wl_loopTM}
\leanok
\uses{Turing.MultiTapeTM, Turing.RecSym, Turing.phaseHostTM, Turing.wl_Ctrl, Turing.wl_red, Turing.wl_rows}
Given a natural number $K$, types $Q,T,P_t,P_s$, Turing machines $R_t$ with one work
tape and $R_s$ with $K+1$ work tapes, and states $t_e,t_y,t_n\in P_t$ and
$s_e,s_y,s_n\in P_s$, assume that equality is decidable on $P_t$ and $P_s$. The
definition constructs a Turing machine with $K+1$ work tapes over the record alphabet
and control states in the dependent sum of the four phases (anchor, test, step, and
stop); its initial control is the anchor phase with unit value, and its transitions use
the phase-specific machines and phase changes determined by the given states.
\end{definition}

\begin{definition}[Configuration with a prescribed state]
\lean{Turing.wl_at}
\label{Turing.wl_at}
\leanok
\uses{Turing.Cfg}
Let $d$ be a configuration with $k$ work tapes over an alphabet $\Gamma$, state set $P$,
and fixed input $x$, and let $q$ belong to a type $Q$. The resulting configuration has
state $q$ and retains $d$’s input-head position, work-tape contents, work-tape head
positions, and output.
\end{definition}

\begin{definition}[Patch the first work tape]
\lean{Turing.wl_patch}
\label{Turing.wl_patch}
\leanok
\uses{Turing.Cfg, Turing.fieldsTape, Turing.joinFields}
Let $d$ be a configuration with $K+1$ work tapes over an alphabet $\Gamma$, state set
$P$, and fixed input $x$; let $q\in P$ and let $f_1,\ldots,f_m$ be finite fields over
$\Gamma$. The resulting configuration has state $q$ and otherwise retains $d$’s
components, except that on work tape $0$ the cells from the head position through the
next $L-1$ cells are set to the concatenated fields separated by blanks, where $L$ is
the length of that concatenation; all other tape cells retain their previous contents.
\end{definition}

\begin{theorem}[A recorded patch changes only the state]
\lean{Turing.wl_patch_self}
\label{Turing.wl_patch_self}
\leanok
\uses{Turing.Cfg, Turing.RecordAt, Turing.wl_at, Turing.wl_patch}
\difficulty{4}
Let $K\in\mathbb{N}$, let $\Gamma$ and $P$ be types, let $x$ be a word over $\Gamma$,
and let $d$ be a configuration with $K+1$ work tapes and states in $P$. If tape $0$ of
$d$ records a field list $f_s$ at every integer offset from $-1$ through the length of
its joined field word, then patching $d$ with state $q$ and fields $f_s$ gives the
configuration obtained from $d$ by changing its state to $q$ and leaving all other data
unchanged.
\end{theorem}

\begin{theorem}[Idempotence of an equal-length tape patch]
\lean{Turing.wl_patch_patch}
\label{Turing.wl_patch_patch}
\leanok
\uses{Turing.Cfg, Turing.joinFields, Turing.wl_patch}
\difficulty{3}
Let $K\in\mathbb{N}$, let $\Gamma$ and $P$ be types, let $x$ be a word over $\Gamma$,
and let $d$ be a configuration with $K+1$ work tapes and state set $P$. Given states
$q,q'\in P$ and lists of fields $f_s,g_s$ over $\Gamma$, suppose their joined words have
equal lengths. Patching $d$ with $f_s$ and state $q$, then patching the result with
$g_s$ and state $q'$, gives the same configuration as patching the original $d$ with
$g_s$ and state $q'$. Each patch sets the state and overwrites work tape $0$ on the
interval beginning at its current head position whose length is the length of the joined
fields, leaving all other configuration data unchanged.
\end{theorem}

\begin{theorem}[A patch records its field list]
\lean{Turing.wl_patch_record}
\label{Turing.wl_patch_record}
\leanok
\uses{Turing.Cfg, Turing.RecordAt, Turing.fieldsTape, Turing.joinFields, Turing.wl_patch}
\difficulty{5}
Let $K\in\mathbb{N}$, let $\Gamma$ and $P$ be types, let $x$ be a word over $\Gamma$,
and let $d$ be a configuration with $K$ work tapes and state set $P$. Suppose tape $r$
of $d$ records the field list $f_s$ at every integer offset $p$ from its head position
with $-1\le p\le L$, where $L$ is the length of the joined field word for $f_s$. If
$g_s$ has the same joined-word length as $f_s$, then patching $d$ with state $q$ and
fields $g_s$ makes tape $0$ record $g_s$ at every such offset from its head position.
\end{theorem}

\begin{theorem}[A phase segment in the hosted machine]
\lean{Turing.wl_phase_segment}
\label{Turing.wl_phase_segment}
\leanok
\uses{Turing.Cfg, Turing.Cfg.mapState, Turing.MultiTapeTM, Turing.MultiTapeTM.runFrom, Turing.RecSym, Turing.Segment, Turing.WindowExit, Turing.phaseHostTM_run, Turing.wl_Ctrl, Turing.wl_Phase, Turing.wl_loopTM, Turing.wl_red, Turing.wl_rows}
\difficulty{4}
Let $K\in\mathbb{N}$, let $Q,T,P_t,P_s$ be types with decidable equality on $P_t,P_s$,
and let $R_t$ and $R_s$ be machines on respectively one and $K+1$ work tapes over the
record alphabet. Consider a phase $i$ other than the stop phase, configurations $a,b$
for the row machine of phase $i$, a time $t$, and a phase value $v$ such that running
that row machine from $a$ for $t$ steps reaches $b$, whose state is $v$. Suppose the
phase transition sends $(i,v)$ to a hosted control state $q$, and no state encountered
strictly before time $t$ admits a phase transition. If a predicate $B$ holds of all
work-head position vectors throughout the row run, including at time $t$, then the
hosted machine reaches the state obtained by embedding $b$ in phase $i$ and replacing
its phase control by $q$ after $t+1$ steps. Before that final step it visits no state in
the stop phase, and $B$ holds of its work-head position vector at every step through
$t+1$.
\end{theorem}

\begin{theorem}[Fixed configurations at stop states]
\lean{Turing.wl_stop_fixed}
\label{Turing.wl_stop_fixed}
\leanok
\uses{Turing.Action.apply, Turing.Action.mapState, Turing.Cfg, Turing.MultiTapeTM, Turing.MultiTapeTM.runFrom, Turing.MultiTapeTM.step, Turing.RecSym, Turing.WindowExit, Turing.phaseHostTM, Turing.wl_Ctrl, Turing.wl_loopTM, Turing.wl_red, Turing.wl_rows}
\difficulty{3}
Let $K$ be a natural number, let $Q,T,P_t,P_s$ be types, and let $R_t,R_s$ be Turing
machines with respectively one and $K+1$ work tapes over record symbols. Choose test
states $t_e,t_y,t_n$ and step states $s_e,s_y,s_n$. For any input and configuration $d$
of the loop machine, any exit label $e$, and any natural number $t$, if $d$ is in stop
state $e$, then after $t$ steps the loop machine is still at exactly $d$.
\end{theorem}

\begin{theorem}[One round of the window loop]
\lean{Turing.wl_round}
\label{Turing.wl_round}
\leanok
\uses{Turing.Cfg, Turing.Cfg.InWindow, Turing.Cfg.mapState, Turing.Cfg.select, Turing.MultiTapeTM, Turing.MultiTapeTM.outputSymbol, Turing.MultiTapeTM.runFrom, Turing.MultiTapeTM.step, Turing.RecSym, Turing.RecSym.digitEmb, Turing.Segment, Turing.Segment.trans, Turing.WindowExit, Turing.cfgRecord, Turing.embedEmitCfg, Turing.embedEmitCfg_eq, Turing.embedEmitCfg_select, Turing.embedEmitCfg_selected_pos, Turing.embedEmitCfg_unselected, Turing.embedEmitTM, Turing.embedEmitTM_runFrom, Turing.fieldsTape, Turing.joinFields, Turing.joinFields_cfgRecord_length, Turing.recordLength, Turing.wl_Ctrl, Turing.wl_at, Turing.wl_loopTM, Turing.wl_patch, Turing.wl_phase_segment, Turing.wl_recordTape, Turing.wl_red, Turing.wl_rows}
\difficulty{7}
Let $k,K$ be nonnegative integers, let $Q,T,P_t,P_s$ be types, let $M$ be a $k$-tape
Turing machine on Boolean symbols, and let $R_t$ and $R_s$ be Turing machines on record
symbols with respectively one and $K+1$ work tapes and state types $P_t$ and $P_s$.
Choose states $t_e,t_y,t_n\in P_t$ and $s_e,s_y,s_n\in P_s$ with $t_y\ne t_n$ and
$s_y\ne s_n$, and natural numbers $C_t,C_s$. Given a Boolean input word $x$, a window
radius $s$, a configuration $c$ of $M$ on $x$, a tag $\tau\in T$, and a configuration
$d$ of the loop machine on the encoded input whose state is the anchor state, suppose
the test machine, started at state $t_e$ on the record tape selected from $d$, has a run
of length at most $C_t(L+1)$ that avoids states $t_y,t_n$ at every earlier time, moves
its head by at most $C_t(L+1)$ from the initial tape-$0$ head position, and ends in
state $t_y$ if $c$ is halted and in state $t_n$ otherwise. Here $L$ is the length of the
joined fields of the record of $(s,c,\tau)$. Suppose also that the step machine, started
from $d$ at state $s_e$, has a run of length at most $C_s(L+|x|+1)$ that avoids states
$s_y,s_n$ at every earlier time and moves each work head by at most $C_s(L+1)$ from its
initial position. If the next configuration $M(c)$ lies within radius $s$, this run ends
in state $s_y$ with tape $0$ replaced on the record interval by the record of $M(c)$ and
the updated tag; if it does not, the run ends in state $s_n$ with the configuration
otherwise unchanged. Then the loop machine has a segment of length $t$ with $1\le
t\le(C_t+C_s+3)(\operatorname{recordLength}(k,|x|,s)+|x|+1)$, avoiding its stop states
at all earlier times and keeping every work-head position within
$(C_t+C_s+3)(\operatorname{recordLength}(k,|x|,s)+1)$ of its initial position through
time $t$. Its endpoint is the halted stop state if $c$ is halted; otherwise it is the
anchor state with the updated record when $M(c)$ lies within the window, and the
left-exit stop state when it does not.
\end{theorem}

\begin{theorem}[Simulating a bounded window run]
\lean{Turing.exists_windowLoopTM}
\label{Turing.exists_windowLoopTM}
\leanok
\uses{Turing.Cfg, Turing.Cfg.InWindow, Turing.Cfg.select, Turing.FinTM, Turing.MultiTapeTM, Turing.MultiTapeTM.outputSymbol, Turing.MultiTapeTM.output_prefix, Turing.MultiTapeTM.runFrom, Turing.MultiTapeTM.runFrom_add, Turing.MultiTapeTM.step, Turing.RecSym, Turing.RecSym.digitEmb, Turing.RecordAt, Turing.Segment, Turing.Segment.trans, Turing.Segment.zero, Turing.WindowExit, Turing.cfgRecord, Turing.exists_recordStepTM, Turing.exists_recordTestTM, Turing.fieldsTape, Turing.joinFields_cfgRecord_length, Turing.recordLength, Turing.windowRun, Turing.wl_Ctrl, Turing.wl_at, Turing.wl_loopTM, Turing.wl_patch, Turing.wl_patch_patch, Turing.wl_patch_record, Turing.wl_patch_self, Turing.wl_recordTape, Turing.wl_round, Turing.wl_stop_fixed}
\difficulty{7}
Let $M$ be a finite-state deterministic Boolean machine, let $T$ be a finite set, and
let $u:T\times\mathrm{Option}(\{0,1\})\to T$ satisfy $u(\tau,\mathrm{none})=\tau$. There
exist a natural number $K$, a finite-state machine $L$ with $K+1$ work tapes, an entry
state, an injective map assigning distinct stop states to the three window-exit
outcomes, and a constant $C$. For every input $x$, window radius $s$, initial
configuration $c$ whose work heads lie in $[-s,s]$ and whose tapes are blank outside
that interval, and tag $\tau$, a correctly encoded record of $(c,\tau)$ on tape 0
together with blank remaining tapes allows $L$ to simulate the window-limited run of $M$
for any number $N$ of steps. It produces the record of the resulting configuration and
updated tag, preserving the initial output tape, and runs in at most $C N(R+|x|+1)$
steps, where $R$ is the prescribed record length; if the window-limited run exhausts all
$N$ steps, then $N$ is no greater than the simulation time, and the machine's status is
the entry state. If the run instead halts or its next step leaves the window, its status
is the corresponding distinct stop state, which it does not enter earlier. At every
time, each work-tape head moves by at most $C(R+1)$ from its initial location.
\end{theorem}

\begin{definition}[Boolean-indexed runner controls]
\lean{Turing.wr_Ctrl}
\label{Turing.wr_Ctrl}
\leanok
Given a type $Q$, define a type depending on a Boolean $b$ to be $Q$ when $b$ is true
and the singleton type when $b$ is false.
\end{definition}

\begin{definition}[Timeout-controlled runner machine]
\lean{Turing.wr_runnerTM}
\label{Turing.wr_runnerTM}
\leanok
\uses{Turing.MultiTapeTM, Turing.phaseHostTM, Turing.wr_Ctrl}
Let $H$ be a $k$-work-tape Turing machine over alphabet $\Gamma$ with state set $Q$, and
let $S$ be a finite type. Given a function assigning a state in $Q$ to each optional
element of $S$, the definition forms a phase-host machine that starts in the true phase
at $H$’s initial state and uses $H$ in that phase. Whenever the current state equals the
assigned state for some element of $S$, the phase transition switches to the false
phase; in the false phase the machine has a singleton state and makes stationary,
nonwriting self-loops.
\end{definition}

\begin{definition}[Window-exit runner controls]
\lean{Turing.wr_exit}
\label{Turing.wr_exit}
\leanok
\uses{Turing.WindowExit, Turing.wr_Ctrl}
Let $P,Q$ be types, let $\mathrm{finish}:\mathrm{Option}(P)\to Q$, and let
$\mathrm{stop}:\mathrm{WindowExit}\to P$. Map the halted and left exits to the true
phase, with control value $\mathrm{finish}$ applied to the corresponding stopped value;
map the exhausted exit to the false phase with the singleton control value.
\end{definition}

\begin{theorem}[Clock-tape configuration decomposition]
\lean{Turing.wr_clockCfg_eta}
\label{Turing.wr_clockCfg_eta}
\leanok
\uses{Turing.Cfg, Turing.Cfg.mapState, Turing.clockCfg, Turing.wr_Ctrl}
\difficulty{4}
Let $K$ be a natural number, let $\Gamma,P,Q$ be types, and let $d$ be a configuration
with $K+2$ work tapes whose state is $(\mathrm{true},q)$ for some $q\in Q$. Then $d$ is
obtained by mapping the state of a configuration with $K+1$ work tapes to
$(\mathrm{true},\cdot)$, where that smaller configuration has state $p$ for some $p\in
P$, retains the input position and output of $d$, and uses the first $K+1$ tapes of $d$;
the additional last tape and its head position are the last tape and head position of
$d$.
\end{theorem}

\begin{theorem}[Agreement before a timeout]
\lean{Turing.wr_prefix}
\label{Turing.wr_prefix}
\leanok
\uses{Turing.Cfg, Turing.Cfg.mapState, Turing.MultiTapeTM, Turing.MultiTapeTM.runFrom, Turing.MultiTapeTM.runFrom_mapState_of_agreeOn, Turing.phaseHostTM, Turing.wr_Ctrl, Turing.wr_runnerTM}
\difficulty{4}
Let $H$ be a Turing machine with state set $Q$, let $S$ be a finite type, and let
$\mathrm{timeout}:\mathrm{Option}(S)\to Q$. If a run of $H$ from configuration $a$
avoids every state $\mathrm{timeout}(r)$ for all times strictly before $t$, then for
every $u\le t$, the runner machine started from $a$ with its active-mode flag set agrees
with the run of $H$ through time $u$, including the configuration state with that flag
set.
\end{theorem}

\begin{theorem}[Timeout dispatch after the first hit]
\lean{Turing.wr_timeout_dispatch}
\label{Turing.wr_timeout_dispatch}
\leanok
\uses{Turing.Cfg, Turing.Cfg.mapState, Turing.MultiTapeTM, Turing.MultiTapeTM.runFrom, Turing.phaseHostTM_run, Turing.wr_Ctrl, Turing.wr_runnerTM}
\difficulty{3}
Let $H$ be a $k$-tape Turing machine with state set $Q$, let $S$ be a finite type, and
let $\tau : \mathrm{Option}(S) \to Q$ be a timeout-state map. If the run of $H$ from a
configuration $a$ is in state $\tau(r)$ after $t$ steps for some $r$, and no run at any
earlier time $u<t$ is in any state $\tau(r')$, then the timeout runner, started from $a$
with its state tagged as the running phase, is after $t+1$ steps exactly the
configuration reached by $H$ after $t$ steps, with its state replaced by the halted
runner phase.
\end{theorem}

\begin{theorem}[Space bound from head displacement]
\lean{Turing.wr_space}
\label{Turing.wr_space}
\leanok
\uses{Turing.Cfg, Turing.MultiTapeTM, Turing.MultiTapeTM.runFrom, Turing.MultiTapeTM.spaceUsed, Turing.MultiTapeTM.spaceUsedByTape, Turing.MultiTapeTM.spaceUsedByTape_le_card_Icc}
\difficulty{3}
Let $R$ be a Turing machine with $k$ work tapes, and let $d$ be an initial
configuration. If, at every time $u\le t$, every work-tape head lies within distance $r$
of its initial position, then the total number of distinct work-tape cells visited
through time $t$ is at most $k(2r+1)$.
\end{theorem}

\begin{theorem}[Finite-window simulation by a Turing machine]
\lean{Turing.exists_windowRunnerTM}
\label{Turing.exists_windowRunnerTM}
\leanok
\uses{Turing.Cfg, Turing.Cfg.InWindow, Turing.Cfg.Stopped, Turing.Cfg.mapState, Turing.FinTM, Turing.FinTM.bufferTape, Turing.MultiTapeTM, Turing.MultiTapeTM.runFrom, Turing.MultiTapeTM.runFrom_state_ne_none_of_le, Turing.MultiTapeTM.spaceUsed, Turing.RecSym, Turing.RecSym.digitEmb, Turing.RecordAt, Turing.WindowExit, Turing.cfgRecord, Turing.clockCfg, Turing.exists_stepClockTMVia, Turing.exists_windowLoopTM, Turing.fieldsTape, Turing.recordLength, Turing.windowRun, Turing.windowRun_stop_unique, Turing.wordPatch, Turing.wr_Ctrl, Turing.wr_clockCfg_eta, Turing.wr_exit, Turing.wr_prefix, Turing.wr_runnerTM, Turing.wr_space, Turing.wr_timeout_dispatch}
\difficulty{7}
Let $M$ be a finite-state Turing machine with Boolean work-tape symbols, and let $T$ be
a finite set with decidable equality. Suppose an update rule
$u:T\times(\{0,1\}\cup\{\mathrm{blank}\})\to T$ satisfies $u(\tau,\mathrm{blank})=\tau$
for every $\tau\in T$. There exist a natural number $K$, a finite state set $P$ with
decidable equality, a machine $R$ with $K+2$ work tapes over an alphabet encoding
Boolean digits, work-tape cells, states, and tags, an entry state $p_{\mathrm{in}}$, an
injective map $e$ from the three outcomes halted, left, and exhausted to states of $P$,
and a natural number $C$, such that every $e(o)$ differs from $p_{\mathrm{in}}$. For
every Boolean input $x$, window size $s$, configuration $c$ of $M$ whose work-tape heads
lie within distance $s$ of zero and whose tapes are blank outside that window, tag
$\tau\in T$, Boolean word $w$, and initial configuration $d$ of $R$ on the digit
encoding of $x$, assume that $d$ is in state $p_{\mathrm{in}}$, tape 0 records the
encoded configuration $c$ and tag $\tau$, the last tape agrees with the digit encoding
of $w$ at every offset from $-1$ through $|w|$, and every other tape is blank. Then
there are a time $t$ and an outcome $o$ such that $t\le 5V(w)+2|w|+3$, where $V(w)$ is
the base-two value of the digits of $w$ in their given order; at time $t$, $R$ is in
state $e(o)$, and it was in no state $e(o')$ at any earlier time. Up to time $t$, each
work-tape head moves by at most $C(L+|w|+1)$ from its initial position, the output is
unchanged, and the total number of work-tape cells visited is at most
$(K+2)(2C(L+|w|+1)+1)$, where $L=k(2s+1)+\operatorname{posWidth}(|x|)+k+4$ and $k$ is
the number of work tapes of $M$. For every natural number $N$, if the windowed run of
$M$ for $N$ steps has an outcome other than exhausted and $CN(L+|x|+1)\le V(w)$, then
$o$ is that outcome. If $o$ is not exhausted, there is some $N$ for which the windowed
run has outcome $o$; at time $t$, tape 0 records its resulting configuration and updated
tag, and the last tape agrees, at every integer cell, with the original tape patched by
a Boolean word of length $|w|$ at its initial head position; all other parts of $R$'s
configuration are unchanged.
\end{theorem}
